\PassOptionsToPackage{table,svgnames,dvipsnames}{xcolor}

\usepackage[utf8]{inputenc}
\usepackage[T1]{fontenc}
\usepackage[sc]{mathpazo}

\usepackage[english]{babel}
\usepackage[autostyle]{csquotes}
\usepackage[sorting=none, autolang=other, backend=biber, style=numeric-comp, citestyle=numeric-comp,
	url=true, isbn=true, minnames=2, giveninits=true]
	{biblatex} %minnames can also be set to 1
\usepackage{textalpha}
\usepackage{graphicx}
\usepackage{scrhack} % necessary for listings package
\usepackage{listings}
\usepackage{lstautogobble}
\usepackage{tikz}
\usepackage[compat=1.1.0]{tikz-feynman}
%\tikzfeynmanset{warn luatex=false}
\usepackage{pgfplots}
\usepackage{tabularx}
\usepackage{pgfplotstable}
\usepackage{booktabs} % for better looking table creations, but bad with vertical lines by design (package creator despises vertical lines)
\usepackage[final]{microtype}
\usepackage{caption}
\usepackage[hidelinks]{hyperref} % hidelinks removes colored boxes around references and links
\usepackage{ifthen} % for comparison of the current language and changing of the thesis layout
\usepackage{pdftexcmds} % string compare to work with all engines
\usepackage{paralist} % for condensed enumerations or lists
%\usepackage{subfig} % for having figures side by side
%\usepackage{siunitx} % for physical accurate units and other numerical presentations
\usepackage{multirow} % makes it possible to have bigger cells over multiple rows in a table
\usepackage{array} % different options for table cell orientation
\usepackage{makecell} % allows nice manual configuration of cells with linebreaks in \thead and \makecell with alignments
\usepackage{pdfpages} % for including multiple pages of pdfs
\usepackage{adjustbox} % can center content wider than the \textwidth
\usepackage{tablefootnote} % for footnotes in tables as \tablefootnote
\usepackage[bottom]{footmisc}
%\interfootnotelinepenalty=10000
\usepackage{threeparttable} % another way to add footnotes as \tablenotes with \item [x] <your footnote> after setting \tnote{x} 
\usepackage{xcolor}
\usepackage{wrapfig}
\usepackage{slashed}
%----------Christoph ANFANG------------------------
\usepackage[version=4]{mhchem} 	% chemistry stuff % set version, otherwise it complains
\usepackage[separate-uncertainty=true, per-mode=symbol]{siunitx}
\usepackage{subcaption} 	% use multiple figures in one figure environment; defines the subfigure environment and replaces the now outdated subfigure package
\usepackage{float}
\usepackage{amsmath,amssymb} 		% math stuff
\usepackage{mathtools}
%\usepackage{mathabx} % enables \corresponds but makes \propto shitty
\usepackage[section]{placeins} % make floats not float into the next section
\usepackage{verbatim} 	% allows to use comment environment
\usepackage{eurosym} 	% the € symbol
% begin quote
\newlength\longest

% chapter quotes
\setkomafont{dictumtext}{\itshape\small}
\setkomafont{dictumauthor}{\normalfont}
\renewcommand*\dictumwidth{\linewidth}
\renewcommand*\dictumauthorformat[1]{--- #1}
\renewcommand*\dictumrule{}
% set depth of the table of content to 2, i.e. also show subsection
\setcounter{tocdepth}{2}
% add costum units to siunitx
\DeclareSIUnit\year{yr}
\DeclareSIUnit\yr{yr}
\DeclareSIUnit\ton{t}
\DeclareSIUnit\t{t}
\DeclareSIUnit\h{h}
\DeclareSIUnit\ppm{ppm}
\DeclareSIUnit\ppmv{ppm(v)}
\DeclareSIUnit\ppmm{ppm(m)}
\DeclareSIUnit\Bq{\becquerel}
\DeclareSIUnit{\sieuro}{\mbox{\euro}}
\DeclareSIUnit{\photoelectron}{pe}

% define some macros
\newcommand{\nldb}{$0\nu\beta\beta$}
\newcommand{\nadb}{$2\nu\beta\beta$}
\newcommand{\code}[1]{\texttt{#1}}
\newcommand{\GE}{\textsc{Gerda}}
\newcommand{\MA}{\textsc{Majorana Demonstrator}}
\newcommand{\LE}{\textsc{Legend}}
\newcommand{\legend}{\textsc{Legend}}
\newcommand{\ltwo}{\textsc{Legend-200}}
\newcommand{\lone}{\textsc{Legend-1000}}
\newcommand{\idefix}{\textsc{Idefix}}
\newcommand{\scarf}{\textsc{Scarf}}
\newcommand{\llama}{\textsc{Llama}}
\newcommand{\llamatwo}{\textsc{Llama-ii}}
\newcommand{\llars}{\textsc{LlArS}}
\newcommand{\Qbb}{$Q_{\beta\beta}$}
\newcommand{\Ge}{$^{76}$Ge}
\newcommand{\bra}[1]{\left<#1\right|}			% der Bra-Vektor <#1|
\newcommand{\ket}[1]{\left|#1\right>}			% der Ket-Vektor |#1>
\newcommand{\braket}[2]{\left<\left.#1\vphantom{#2}\right|#2\right>}		% das Skalarprodukt <#1|#2>
\newcommand{\e}{\mathrm{e}}					% die Eulersche Zahl e
%\newcommand{\O}{\mathcal{O}\left(#1\right)}
\newcommand{\ida}{\vcentcolon=} % is defined as
\newcommand{\corresponds}{\overset{\scriptscriptstyle\wedge}{=}}

\newcommand{\cn}{$^{\text{\textcolor{blue}{citation needed}}}$}
\newcommand{\ms}{$^\text{\textcolor{green}{make sure}}$}
\newcommand{\fn}{$^\text{\textcolor{red}{fill number}}$}
%\newcommand{\ac}{$^\text{\textcolor{magenta}{add content}}$}

% redefinde \autoref{}, a feature of hyperref, such that sections, etc. are capitalized
\addto\extrasenglish{\def\chapterautorefname{Chapter}}
\addto\extrasenglish{\def\sectionautorefname{Section}}
\addto\extrasenglish{\def\subsectionautorefname{Section}}
\addto\extrasenglish{\def\subsubsectionautorefname{Section}}
\addto\extrasenglish{\def\paragraphautorefname{Paragraph}}
%---------Christoph ENDE--------------------------
% https://tex.stackexchange.com/questions/42619/x-mark-to-match-checkmark
\usepackage{amssymb}% http://ctan.org/pkg/amssymb
\usepackage{pifont}% http://ctan.org/pkg/pifont
\newcommand{\cmark}{\ding{51}}%
\newcommand{\xmark}{\ding{55}}%


\usepackage[acronym,xindy,toc]{glossaries}
\usepackage{acronym} %include "acronym" if glossary and acronym should be separated
\addbibresource{bibliography.bib}
%\bibliography{bibliography}
%--------------------------------------------------
\AtEveryBibitem{
	\clearfield{pages} 		% do not show page number
	\clearfield{pagetotal} 	% also for books
}
%--------------------------------------------------
\linespread{1.05} % adjust line spread for mathpazo font

% Add table of contents to PDF bookmarks
\BeforeTOCHead[toc]{{\cleardoublepage\pdfbookmark[0]{\contentsname}{toc}}}

% Define TUM corporate design colors
% Taken from http://portal.mytum.de/corporatedesign/index_print/vorlagen/index_farben
\definecolor{TUMBlue}{HTML}{0065BD}
\definecolor{TUMSecondaryBlue}{HTML}{005293}
\definecolor{TUMSecondaryBlue2}{HTML}{003359}
\definecolor{TUMBlack}{HTML}{000000}
\definecolor{TUMWhite}{HTML}{FFFFFF}
\definecolor{TUMDarkGray}{HTML}{333333}
\definecolor{TUMGray}{HTML}{808080}
\definecolor{TUMLightGray}{HTML}{CCCCC6}
\definecolor{TUMAccentGray}{HTML}{DAD7CB}
\definecolor{TUMAccentOrange}{HTML}{E37222}
\definecolor{TUMAccentGreen}{HTML}{A2AD00}
\definecolor{TUMAccentLightBlue}{HTML}{98C6EA}
\definecolor{TUMAccentBlue}{HTML}{64A0C8}

% Settings for pgfplots
\pgfplotsset{compat=newest}
\pgfplotsset{
  % For available color names, see http://www.latextemplates.com/svgnames-colors
  cycle list={TUMBlue\\TUMAccentOrange\\TUMAccentGreen\\TUMSecondaryBlue2\\TUMDarkGray\\},
}

% Settings for lstlistings

% Use this for basic highlighting
\lstset{%
  basicstyle=\ttfamily,
  columns=fullflexible,
  autogobble,
  keywordstyle=\bfseries\color{TUMBlue},
  stringstyle=\color{TUMAccentGreen}
}

% use this for C# highlighting
% %\setmonofont{Consolas} %to be used with XeLaTeX or LuaLaTeX
% \definecolor{bluekeywords}{rgb}{0,0,1}
% \definecolor{greencomments}{rgb}{0,0.5,0}
% \definecolor{redstrings}{rgb}{0.64,0.08,0.08}
% \definecolor{xmlcomments}{rgb}{0.5,0.5,0.5}
% \definecolor{types}{rgb}{0.17,0.57,0.68}

% \lstset{language=[Sharp]C,
% captionpos=b,
% %numbers=left, % numbering
% %numberstyle=\tiny, % small row numbers
% frame=lines, % above and underneath of listings is a line
% showspaces=false,
% showtabs=false,
% breaklines=true,
% showstringspaces=false,
% breakatwhitespace=true,
% escapeinside={(*@}{@*)},
% commentstyle=\color{greencomments},
% morekeywords={partial, var, value, get, set},
% keywordstyle=\color{bluekeywords},
% stringstyle=\color{redstrings},
% basicstyle=\ttfamily\small,
% }

% Settings for search order of pictures
\graphicspath{
    {logos/}
    {figures/}
}

% Set up hyphenation rules for the language package when mistakes happen
\babelhyphenation[english]{
an-oth-er
ex-am-ple
}

\newcommand{\todo}[1]{{\bfseries{\scshape{\color{TUMAccentOrange}[(TODO: #1)]}}}} % for multiple paragraphs
\newcommand{\done}[1]{{\itshape{\scshape{\color{TUMAccentBlue}[(Done: #1)]}}}} % for multiple paragraphs
% for error handling of intended behavior in your latex documents.

\newcommand{\tabitem}{~~\llap{\textbullet}~~}

\newcolumntype{P}[1]{>{\centering\arraybackslash}p{#1}} % for horizontal alignment with limited column width
\newcolumntype{M}[1]{>{\centering\arraybackslash}m{#1}} % for horizontal and vertical alignment with limited column width
\newcolumntype{L}[1]{>{\raggedright\arraybackslash}m{#1}} % for vertical alignment left with limited column width
\newcolumntype{R}[1]{>{\raggedleft\arraybackslash}m{#1}} % for vertical alignment right with limited column width